\documentclass[stu,letterpaper,12pt,donotrepeattitle,floatsintext]{apa7}
% Preámbulo común de los documentos de Nexus: Normas APA 7.a edición (trabajo de estudiante).
% Se carga con % Preámbulo común de los documentos de Nexus: Normas APA 7.a edición (trabajo de estudiante).
% Se carga con % Preámbulo común de los documentos de Nexus: Normas APA 7.a edición (trabajo de estudiante).
% Se carga con \input{../comun/preambulo} justo después de \documentclass.
\usepackage[spanish,es-noshorthands,es-tabla]{babel}
\usepackage[T1]{fontenc}
\usepackage{newtxtext,newtxmath}
\usepackage{csquotes}
\usepackage[style=apa,backend=biber]{biblatex}
\addbibresource{../comun/referencias.bib}
\usepackage{graphicx}
\usepackage{booktabs}
\usepackage{tabularx}
\usepackage{array}
\usepackage{enumitem}
\usepackage{xurl}
\usepackage{tikz}
\usetikzlibrary{positioning,arrows.meta,shapes.misc}
\hypersetup{hidelinks}
\setlist{itemsep=2pt,topsep=4pt,parsep=0pt}
\setlength{\emergencystretch}{3em}

% Datos de la portada
\authorsnames{Neyder Sebastián Orozco Villamil, Deisy Lorena Jiménez Martínez}
\authorsaffiliations{Universidad Central}
\course{Modelado 3D y Video Juegos (43391312)}
\professor{Mg. Carlos Iván Pinzón Romero}
\duedate{18 de noviembre de 2026}

% Atajos de escritura
\newcommand{\cod}[1]{\texttt{#1}}            % archivos, scripts y código
\newcommand{\ruta}[1]{\textit{#1}}           % rutas de menú: Archivo > Guardar
\newcommand{\tecla}[1]{\textsc{#1}}          % teclas

% Columnas de tabla con texto justificado a la izquierda
\newcolumntype{L}{>{\raggedright\arraybackslash}X}
\newcolumntype{P}[1]{>{\raggedright\arraybackslash}p{#1}}

% Ubicación de figuras y tablas: cerca del texto (los manuales usan H: exactamente en su lugar)
\usepackage{float}
\floatplacement{figure}{!htbp}
\floatplacement{table}{!htbp}
% Nota de figuras y tablas en español y alineada a la izquierda (APA)
\AtBeginDocument{\renewcommand{\notelabel}{Nota}}
\newcommand{\NotaFig}[1]{\if\relax\detokenize{#1}\relax\else\par\raggedright\figurenote{#1}\fi}

% Figura APA: número y título arriba, imagen y nota opcional abajo
%   \Figura[ancho]{archivo}{etiqueta}{título}{nota}
\newcommand{\Figura}[5][0.82\linewidth]{%
  \begin{figure}
    \caption{#4}\label{#3}
    \centering\includegraphics[width=#1]{#2}
    \NotaFig{#5}
  \end{figure}}

% Dos imágenes lado a lado en una misma figura
%   \FiguraDoble{archivo a}{archivo b}{etiqueta}{título}{nota}
\newcommand{\FiguraDoble}[5]{%
  \begin{figure}
    \caption{#4}\label{#3}
    \centering
    \includegraphics[width=0.49\linewidth]{#1}\hfill\includegraphics[width=0.49\linewidth]{#2}
    \NotaFig{#5}
  \end{figure}}

% Tabla APA: número y título arriba; contenido a espacio sencillo y en letra un poco menor
\newenvironment{TablaAPA}[2]{%
  \begin{table}\caption{#2}\label{#1}\centering\small\linespread{1}\selectfont}
  {\end{table}}

% Términos de las referencias como en el manual de normas APA adjunto (Menéndez Echavarría, 2020)
\DefineBibliographyStrings{spanish}{version = {Versión}, nodate = {s.\,f\adddot}}
 justo después de \documentclass.
\usepackage[spanish,es-noshorthands,es-tabla]{babel}
\usepackage[T1]{fontenc}
\usepackage{newtxtext,newtxmath}
\usepackage{csquotes}
\usepackage[style=apa,backend=biber]{biblatex}
\addbibresource{../comun/referencias.bib}
\usepackage{graphicx}
\usepackage{booktabs}
\usepackage{tabularx}
\usepackage{array}
\usepackage{enumitem}
\usepackage{xurl}
\usepackage{tikz}
\usetikzlibrary{positioning,arrows.meta,shapes.misc}
\hypersetup{hidelinks}
\setlist{itemsep=2pt,topsep=4pt,parsep=0pt}
\setlength{\emergencystretch}{3em}

% Datos de la portada
\authorsnames{Neyder Sebastián Orozco Villamil, Deisy Lorena Jiménez Martínez}
\authorsaffiliations{Universidad Central}
\course{Modelado 3D y Video Juegos (43391312)}
\professor{Mg. Carlos Iván Pinzón Romero}
\duedate{18 de noviembre de 2026}

% Atajos de escritura
\newcommand{\cod}[1]{\texttt{#1}}            % archivos, scripts y código
\newcommand{\ruta}[1]{\textit{#1}}           % rutas de menú: Archivo > Guardar
\newcommand{\tecla}[1]{\textsc{#1}}          % teclas

% Columnas de tabla con texto justificado a la izquierda
\newcolumntype{L}{>{\raggedright\arraybackslash}X}
\newcolumntype{P}[1]{>{\raggedright\arraybackslash}p{#1}}

% Ubicación de figuras y tablas: cerca del texto (los manuales usan H: exactamente en su lugar)
\usepackage{float}
\floatplacement{figure}{!htbp}
\floatplacement{table}{!htbp}
% Nota de figuras y tablas en español y alineada a la izquierda (APA)
\AtBeginDocument{\renewcommand{\notelabel}{Nota}}
\newcommand{\NotaFig}[1]{\if\relax\detokenize{#1}\relax\else\par\raggedright\figurenote{#1}\fi}

% Figura APA: número y título arriba, imagen y nota opcional abajo
%   \Figura[ancho]{archivo}{etiqueta}{título}{nota}
\newcommand{\Figura}[5][0.82\linewidth]{%
  \begin{figure}
    \caption{#4}\label{#3}
    \centering\includegraphics[width=#1]{#2}
    \NotaFig{#5}
  \end{figure}}

% Dos imágenes lado a lado en una misma figura
%   \FiguraDoble{archivo a}{archivo b}{etiqueta}{título}{nota}
\newcommand{\FiguraDoble}[5]{%
  \begin{figure}
    \caption{#4}\label{#3}
    \centering
    \includegraphics[width=0.49\linewidth]{#1}\hfill\includegraphics[width=0.49\linewidth]{#2}
    \NotaFig{#5}
  \end{figure}}

% Tabla APA: número y título arriba; contenido a espacio sencillo y en letra un poco menor
\newenvironment{TablaAPA}[2]{%
  \begin{table}\caption{#2}\label{#1}\centering\small\linespread{1}\selectfont}
  {\end{table}}

% Términos de las referencias como en el manual de normas APA adjunto (Menéndez Echavarría, 2020)
\DefineBibliographyStrings{spanish}{version = {Versión}, nodate = {s.\,f\adddot}}
 justo después de \documentclass.
\usepackage[spanish,es-noshorthands,es-tabla]{babel}
\usepackage[T1]{fontenc}
\usepackage{newtxtext,newtxmath}
\usepackage{csquotes}
\usepackage[style=apa,backend=biber]{biblatex}
\addbibresource{../comun/referencias.bib}
\usepackage{graphicx}
\usepackage{booktabs}
\usepackage{tabularx}
\usepackage{array}
\usepackage{enumitem}
\usepackage{xurl}
\usepackage{tikz}
\usetikzlibrary{positioning,arrows.meta,shapes.misc}
\hypersetup{hidelinks}
\setlist{itemsep=2pt,topsep=4pt,parsep=0pt}
\setlength{\emergencystretch}{3em}

% Datos de la portada
\authorsnames{Neyder Sebastián Orozco Villamil, Deisy Lorena Jiménez Martínez}
\authorsaffiliations{Universidad Central}
\course{Modelado 3D y Video Juegos (43391312)}
\professor{Mg. Carlos Iván Pinzón Romero}
\duedate{18 de noviembre de 2026}

% Atajos de escritura
\newcommand{\cod}[1]{\texttt{#1}}            % archivos, scripts y código
\newcommand{\ruta}[1]{\textit{#1}}           % rutas de menú: Archivo > Guardar
\newcommand{\tecla}[1]{\textsc{#1}}          % teclas

% Columnas de tabla con texto justificado a la izquierda
\newcolumntype{L}{>{\raggedright\arraybackslash}X}
\newcolumntype{P}[1]{>{\raggedright\arraybackslash}p{#1}}

% Ubicación de figuras y tablas: cerca del texto (los manuales usan H: exactamente en su lugar)
\usepackage{float}
\floatplacement{figure}{!htbp}
\floatplacement{table}{!htbp}
% Nota de figuras y tablas en español y alineada a la izquierda (APA)
\AtBeginDocument{\renewcommand{\notelabel}{Nota}}
\newcommand{\NotaFig}[1]{\if\relax\detokenize{#1}\relax\else\par\raggedright\figurenote{#1}\fi}

% Figura APA: número y título arriba, imagen y nota opcional abajo
%   \Figura[ancho]{archivo}{etiqueta}{título}{nota}
\newcommand{\Figura}[5][0.82\linewidth]{%
  \begin{figure}
    \caption{#4}\label{#3}
    \centering\includegraphics[width=#1]{#2}
    \NotaFig{#5}
  \end{figure}}

% Dos imágenes lado a lado en una misma figura
%   \FiguraDoble{archivo a}{archivo b}{etiqueta}{título}{nota}
\newcommand{\FiguraDoble}[5]{%
  \begin{figure}
    \caption{#4}\label{#3}
    \centering
    \includegraphics[width=0.49\linewidth]{#1}\hfill\includegraphics[width=0.49\linewidth]{#2}
    \NotaFig{#5}
  \end{figure}}

% Tabla APA: número y título arriba; contenido a espacio sencillo y en letra un poco menor
\newenvironment{TablaAPA}[2]{%
  \begin{table}\caption{#2}\label{#1}\centering\small\linespread{1}\selectfont}
  {\end{table}}

% Términos de las referencias como en el manual de normas APA adjunto (Menéndez Echavarría, 2020)
\DefineBibliographyStrings{spanish}{version = {Versión}, nodate = {s.\,f\adddot}}

\graphicspath{{figuras/}{../desarrollo/figuras/}}
\floatplacement{figure}{H}
\floatplacement{table}{H}

\title{Manual de instalación del videojuego Nexus para Android (APK)}

% ---------- Ilustraciones esquemáticas de un teléfono (TikZ) ----------
% \Telefono{contenido de la pantalla}: dibuja un teléfono vertical; la pantalla va de
% (0.25,0.45) a (5.15,10.2) y su centro horizontal es x = 2.7.
\newcommand{\Telefono}[1]{%
  \begingroup\linespread{1}\selectfont
  \begin{tikzpicture}[x=1cm,y=1cm,font=\sffamily\small]
    \draw[rounded corners=0.45cm,line width=1.4pt,fill=black!80] (0,0) rectangle (5.4,10.6);
    \fill[white,rounded corners=0.18cm] (0.25,0.45) rectangle (5.15,10.2);
    \node[anchor=west,font=\sffamily\scriptsize] at (0.45,9.95) {12:30};
    \node[anchor=east,font=\sffamily\scriptsize] at (4.95,9.95) {\(\blacktriangledown\) 80\,\%};
    #1
  \end{tikzpicture}\endgroup}
% Cuadro de diálogo centrado en la pantalla: \Dialogo{y}{título}{texto}{botón izq.}{botón der.}
\newcommand{\Dialogo}[5]{%
  \fill[black!30,rounded corners=0.18cm] (0.25,0.45) rectangle (5.15,9.75);
  \node[draw=black!25,fill=white,rounded corners=0.25cm,text width=4.1cm,align=left,
        inner sep=8pt] (d) at (2.7,#1) {\textbf{#2}\\[4pt]{\footnotesize #3}\\[10pt]
        \hfill{\footnotesize\color{blue!70!black}\textbf{#4}\quad\textbf{#5}}};}
\newcommand{\Barra}[1]{%
  \fill[black!6] (0.25,8.9) rectangle (5.15,9.75);
  \node[anchor=west,font=\sffamily\bfseries] at (0.5,9.32) {#1};}

\begin{document}
\maketitle

\section{Introducción}

\textit{Nexus: Protocolo de emergencia} es un videojuego de terror y sigilo en tercera persona,
desarrollado en Unity \parencite{unity2026} como trabajo final de la asignatura Modelado 3D y Video
Juegos. Este manual explica cómo instalar la versión para Android a partir del archivo
\cod{Nexus.apk}, cómo se juega con los controles táctiles y cómo resolver los problemas más
comunes.

Un archivo APK es el paquete con el que se instalan las aplicaciones en Android. Como Nexus no se
publica en la tienda Google Play, se instala directamente desde el archivo; por eso Android pide un
permiso adicional la primera vez, como se explica en el paso~3.

\section{Requisitos del teléfono}

\begin{TablaAPA}{tab:requisitos}{Requisitos del teléfono o la tableta}
\begin{tabularx}{\linewidth}{P{4cm}L}
\toprule
Componente & Requisito \\
\midrule
Versión de Android & Android 8.0 (Oreo) o posterior \\
Procesador & De 64 bits (ARM64), presente en casi todos los teléfonos fabricados desde 2017 \\
Memoria RAM & 3 GB o más \\
Espacio libre & 300 MB (el archivo ocupa 75 MB y el juego instalado, un poco más) \\
Pantalla & Se juega en horizontal; se recomienda una pantalla de 6 pulgadas o más \\
Audio & Audífonos recomendados, para ubicar a los androides por el sonido \\
\bottomrule
\end{tabularx}
\end{TablaAPA}

La versión de Android se consulta en \ruta{Ajustes} \textrightarrow{} \ruta{Acerca del teléfono}
\textrightarrow{} \ruta{Información de software} (el nombre exacto del menú cambia según el
fabricante).

\section{Instalación paso a paso}

Las Figuras~\ref{fig:transferir} a~\ref{fig:instalado} son representaciones esquemáticas de las
pantallas de Android, elaboradas para este manual. El texto y la apariencia exactos dependen de la
versión de Android y de la marca del teléfono, pero los pasos son los mismos.

\subsection{Paso 1: copiar el archivo al teléfono}

Hay dos formas de llevar \cod{Nexus.apk} al teléfono:

\begin{itemize}
  \item \textbf{Con un cable USB:} conecte el teléfono al computador, desbloquéelo y, en la
        notificación de USB, elija \ruta{Transferencia de archivos}. En el Explorador de archivos
        del computador abra el teléfono y copie \cod{Nexus.apk} en la carpeta \cod{Download}
        (Figura~\ref{fig:transferir}).
  \item \textbf{Por Internet:} suba el archivo a una carpeta de Google Drive o envíelo por correo,
        y descárguelo desde el teléfono. Queda guardado en la carpeta \ruta{Descargas}.
\end{itemize}

\begin{figure}
  \caption{Representación esquemática de la copia del APK del computador al teléfono}\label{fig:transferir}
  \centering
  \begin{tikzpicture}[font=\sffamily\small]
    \draw[line width=1.2pt,rounded corners=3pt,fill=black!5] (0,1.2) rectangle (4.4,4);
    \draw[line width=1.2pt] (1.6,1.2) -- (1.4,0.6) -- (3,0.6) -- (2.8,1.2);
    \node[text width=3.8cm,align=center] at (2.2,2.6) {Computador\\[3pt]\ttfamily Nexus.apk\\{\sffamily\footnotesize (75 MB)}};
    \draw[-{Stealth[length=3mm]},line width=1.4pt] (4.7,2.5) -- (7.1,2.5)
      node[midway,above,align=center,font=\sffamily\footnotesize] {Cable USB:\\Transferencia\\de archivos};
    \draw[line width=1.2pt,rounded corners=6pt,fill=black!80] (7.4,0) rectangle (9.6,4.4);
    \fill[white,rounded corners=2pt] (7.55,0.25) rectangle (9.45,4.15);
    \node[text width=1.8cm,align=center,font=\sffamily\scriptsize] at (8.5,2.2) {Teléfono\\[3pt]Carpeta\\\ttfamily Download\\[2pt]\ttfamily Nexus.apk};
  \end{tikzpicture}
  \NotaFig{Ilustración elaborada para este manual.}
\end{figure}

\subsection{Paso 2: abrir el archivo en el teléfono}

Abra la aplicación \ruta{Archivos} (en algunos teléfonos se llama \ruta{Mis archivos} o
\ruta{Gestor de archivos}), entre a \ruta{Descargas} y toque \cod{Nexus.apk}
(Figura~\ref{fig:archivos}). Si descargó el archivo con el navegador, también puede tocar la
notificación de la descarga.

\begin{figure}
  \caption{Representación esquemática del archivo APK en la carpeta Descargas}\label{fig:archivos}
  \centering
  \Telefono{%
    \Barra{Descargas}
    \draw[black!15] (0.45,7.6) -- (4.95,7.6);
    \fill[green!55!black,rounded corners=2pt] (0.55,7.85) rectangle (1.15,8.6);
    \node[anchor=west] at (1.35,8.4) {\textbf{Nexus.apk}};
    \node[anchor=west,font=\sffamily\scriptsize,black!60] at (1.35,8.0) {75,4 MB \textperiodcentered{} Hoy};
    \draw[-{Stealth},line width=1pt,red!70!black] (4.6,7.0) -- (3.4,7.95);
    \node[font=\sffamily\scriptsize,red!70!black] at (4.3,6.7) {Tocar};
  }
  \NotaFig{Ilustración elaborada para este manual.}
\end{figure}

\subsection{Paso 3: permitir la instalación desde esta fuente}

La primera vez, Android bloquea la instalación de aplicaciones que no vienen de Google Play y muestra
un aviso como el de la Figura~\ref{fig:fuentes}, izquierda. Para continuar:

\begin{enumerate}
  \item Toque \ruta{Configuración} en el aviso.
  \item Active \ruta{Permitir de esta fuente} para la aplicación con la que abrió el archivo
        (Figura~\ref{fig:fuentes}, derecha).
  \item Vuelva atrás: aparece la pantalla de instalación.
\end{enumerate}

El mismo permiso se encuentra en \ruta{Ajustes} \textrightarrow{} \ruta{Aplicaciones}
\textrightarrow{} \ruta{Acceso especial} \textrightarrow{} \ruta{Instalar apps desconocidas}. Se
recomienda desactivarlo de nuevo cuando termine la instalación.

\begin{figure}
  \caption{Representación esquemática del aviso de fuentes desconocidas y del permiso}\label{fig:fuentes}
  \centering
  \Telefono{%
    \Dialogo{5.3}{Instalación bloqueada}{Por tu seguridad, no se permite que tu teléfono instale apps desconocidas de esta fuente.}{Cancelar}{Configuración}
  }\hspace{1cm}%
  \Telefono{%
    \fill[black!6] (0.25,8.9) rectangle (5.15,9.75);
    \node[anchor=west,font=\sffamily\bfseries\footnotesize] at (0.45,9.32) {Instalar apps desconocidas};
    \node[anchor=west] at (0.5,8.3) {\textbf{Archivos}};
    \node[anchor=west,font=\sffamily\footnotesize] at (0.5,7.45) {Permitir de esta fuente};
    \fill[blue!60!black,rounded corners=0.2cm] (4.15,7.25) rectangle (4.95,7.65);
    \fill[white] (4.75,7.45) circle (0.15);
    \node[text width=4.3cm,align=left,font=\sffamily\scriptsize,black!60] at (2.7,6.5)
      {Si instalas apps de esta fuente, tu teléfono y tus datos pueden correr riesgos.};
  }
  \NotaFig{Ilustración elaborada para este manual. A la izquierda, el aviso; a la derecha, el permiso ya activado.}
\end{figure}

\subsection{Paso 4: instalar}

En la pantalla de instalación toque \ruta{Instalar} (Figura~\ref{fig:instalar}, izquierda). Como el
juego no proviene de la tienda, Google Play Protect, el sistema de seguridad de Android, puede
analizarlo o advertir que no reconoce al desarrollador \parencite{playprotect}. En ese caso toque
\ruta{Más detalles} y luego \ruta{Instalar de todas formas} (Figura~\ref{fig:instalar}, derecha).
El aviso se debe a que la aplicación no está publicada en Google Play, no a un problema del juego.

\begin{figure}
  \caption{Representación esquemática de la confirmación de instalación y del aviso de Play Protect}\label{fig:instalar}
  \centering
  \Telefono{%
    \Dialogo{5.3}{Nexus}{¿Quieres instalar esta app?}{Cancelar}{Instalar}
  }\hspace{1cm}%
  \Telefono{%
    \Dialogo{5.3}{Bloqueada por Play Protect}{Play Protect no reconoce al desarrollador de esta app.\\[4pt]\color{blue!70!black}Más detalles \(\blacktriangledown\)\\\color{black}\textbf{Instalar de todas formas}}{}{Aceptar}
  }
  \NotaFig{Ilustración elaborada para este manual. El aviso de la derecha puede no aparecer.}
\end{figure}

\subsection{Paso 5: abrir el juego}

Al terminar la instalación toque \ruta{Abrir} (Figura~\ref{fig:instalado}). El juego también queda
en la lista de aplicaciones con el nombre \textit{Nexus}. Se juega con el teléfono en horizontal: si
el teléfono está en vertical, la imagen gira sola. La primera carga tarda unos segundos y luego
aparece la pantalla de título (Figura~\ref{fig:titulo}).

\begin{figure}
  \caption{Representación esquemática de la instalación terminada}\label{fig:instalado}
  \centering
  \Telefono{%
    \Dialogo{5.3}{Nexus}{App instalada.}{Listo}{Abrir}
  }
  \NotaFig{Ilustración elaborada para este manual.}
\end{figure}

\Figura[0.9\linewidth]{118_titulo-logo-universidad.jpg}{fig:titulo}{Pantalla de título del juego}{Captura del juego. En el teléfono se ve igual, en horizontal y a pantalla completa.}

\section{Cómo se juega en el teléfono}

En el menú toque \ruta{Nueva partida}, elija a su técnico con las flechas, toque \ruta{Confirmar}
y luego \ruta{Iniciar misión}. Los diálogos avanzan tocando la pantalla y las cinemáticas se pueden
saltar con un toque.

Durante el juego aparecen los controles táctiles (Figura~\ref{fig:tactiles} y
Tabla~\ref{tab:tactiles}). Solo se muestran cuando el jugador tiene el control: se ocultan en los
menús, los diálogos, el inventario y las cinemáticas. Los botones de apuntar y recargar solo
aparecen con un arma de fuego equipada.

\Figura[0.9\linewidth]{115_controles-tactiles.jpg}{fig:tactiles}{Controles táctiles con la pistola equipada}{Captura del juego con los controles táctiles activados.}

\begin{TablaAPA}{tab:tactiles}{Controles táctiles}
\begin{tabularx}{\linewidth}{P{3.4cm}L}
\toprule
Control & Función \\
\midrule
Joystick (mitad izquierda) & Aparece donde se apoya el pulgar. Arriba y abajo: avanzar y retroceder. A los lados: girar \\
CORRER & Se deja activado; tóquelo otra vez para volver a caminar \\
INVENTARIO & Abre el inventario, el archivo de notas y el mapa \\
INTERACTUAR & Usar terminales, puertas, máquinas y notas, y registrar cuerpos \\
HABILIDAD & Señuelo de red (Neyder) o KILLSWITCH.exe (Deisy); se atenúa mientras se recarga \\
ARMA & Cambia entre las armas que se tengan \\
ATACAR / DISPARAR & Ataque con el cuchillo o disparo. Sin estar apuntando, el técnico apunta solo al androide más cercano y dispara \\
APUNTAR & Se deja activado: el técnico mantiene el arma levantada y gira con el joystick \\
RECARGAR & Recarga el arma de fuego \\
\bottomrule
\end{tabularx}
\end{TablaAPA}

\section{Desinstalación}

Mantenga presionado el ícono de Nexus y elija \ruta{Desinstalar}, o entre a \ruta{Ajustes}
\textrightarrow{} \ruta{Aplicaciones} \textrightarrow{} \ruta{Nexus} \textrightarrow{}
\ruta{Desinstalar}. Al desinstalarlo se borra también la partida guardada.

\section{Solución de problemas}

La Tabla~\ref{tab:problemas} reúne los problemas más frecuentes durante la instalación y el juego,
con su solución.

\begin{TablaAPA}{tab:problemas}{Problemas frecuentes y cómo resolverlos}
\begin{tabularx}{\linewidth}{P{4.6cm}L}
\toprule
Problema & Solución \\
\midrule
No aparece la opción de instalar & Active el permiso de fuentes desconocidas para la aplicación con la que abre el archivo (paso 3) \\
\textquote{App no instalada} & Libere espacio, desinstale una versión anterior de Nexus si la hay, y vuelva a copiar el APK por si se copió incompleto \\
\textquote{La app no es compatible} o error de análisis & El teléfono tiene una versión de Android anterior a la 8.0 o un procesador de 32 bits \\
Aviso de Play Protect & Toque \ruta{Más detalles} y luego \ruta{Instalar de todas formas} (paso 4) \\
El juego va lento o se calienta & Cierre otras aplicaciones, active el modo de alto rendimiento si el teléfono lo tiene y conéctelo al cargador \\
No hay sonido & Suba el volumen multimedia (no el del timbre) y revise que no haya audífonos Bluetooth conectados a otro equipo \\
Los controles táctiles no aparecen & Se ocultan en menús, diálogos, inventario y cinemáticas; cierre el diálogo o el inventario para volver a verlos \\
\bottomrule
\end{tabularx}
\end{TablaAPA}

\printbibliography

\end{document}
